\section{FMR studies of surface and bulk spin-wave modes in a CoFe/PtMn/CoFe multilayer film - \href{http://dx.doi.org/10.1063/1.2839337}{JAP 103, 07B525 (2008)}}
\subsection*{Main points from the Paper}
\begin{enumerate}
    \item Studied the exchange-dominated surface and bulk spin-wave modes in a single period of CoFe/PtMn/CoFe trilayer film.
    \item Multimode spin-wave spectra was observed using the FMR. Also, they have identified a high-order standing spin wave in the spectra and found \textbf{"critical angle"} in the multilayer sample. They have also included the effective surface anisotropy field to describe the results.
    \item They have used FMR to investigate the exchange-dominated surface and bulk spin-wave modes in a single period of CoFe/PtMn/CoFe tri-layer film. Bulk magnetic anisotropy parameters and the g-factor of the magnetic film based on the angle dependence of the main resonance field.
    \item A significant enhancement in the in-plane anisotropy was observed in the range; $\sim57-123$ Oe.
    \item A high-order standin spin wave was identified in our spectra along with a "critical angle" in the trilayer sample. 
    \item An effective surface anisotropy field due to the exchange bias and the dynamic surface spin pinning was also used to describe the results.
    \item The exchange interaction stiffness in the CoFe layers, $J\sim 2.7$ meV.
    \item \underline{Acoustic spin-wave mode:} The acoustic spin wave mode is the lowest-frequency mode where all spins in the material move together and stay in phase across the thickness. e.g. In a thin ferromagnetic film, all magnetic moments tilt and precess together when a microwave field is applied.
    \item \underline{Higher-order standing SW:} Higher-order standing spin wave modes occur when spins form standing patterns across the thickness, with some regions moving opposite to others. e.g. In a thick magnetic film, the magnetization oscillates with one or more nodes across the thickness, similar to standing waves on a stretched string.
    \item These results could be related to the change of surface spin pinning is also said in the paper.
    \item A two-fold symmetry for the out-of-plane geometry is observed when the angular dependence of the resonance filed of the acoustic spin-wave mode is measured.
    \item The excitation of the spin-wave precession is due to the absorption of microwave. 
    \item The effective parameter A changes with changing angles when we rotate the sample, and the magnetic profile changes due to the change of the surface pinning. This results in the change of the spin-wave resonance energy.    
\end{enumerate}
\subsection*{Summary}
In summary, we investigated the basic magnetic excitations in a CoFe/PtMn/CoFe multilayer film using the FMR technique. A higher-order spin-wave mode appears only after the magnetic field is tilted beyond a certain angle. By considering an effective surface anisotropy field, the observed spin-wave resonance spectra can be explained by dynamic pinning of spins at the surfaces.


